% Insert after the weight-six genus-ratio subsection of cm_genus.tex.
% Uses the reduced modular polynomial of newmodular:eq:reduced.

\subsection{A fixed field contains all even-weight genus ratios}

The low-weight identities imply a uniform field-degree theorem.
For a fixed genus character, increasing the weight introduces no
additional algebraic field extension beyond the two seed ratios.

\begin{theorem}[Uniform algebraic degree of genus ratios]
\label{cmgenus:thm:degree}
Under the assumptions of Theorem~\ref{cmgenus:thm:allweights}, put
$K=\mathbb Q(\sqrt e)$ and
\[
 \mathcal K_{D,\varepsilon}=K(R_4(D),R_6(D)).
\]
Then
\[
 [\mathcal K_{D,\varepsilon}:K]\leq6,
 \qquad [\mathcal K_{D,\varepsilon}:\mathbb Q]\leq12,
\]
and every defined ratio $R_w(D)$ with even $w\geq4$ belongs to this
same field. The degree bound is independent of $w$ and of the imaginary
class number $h(D)$.
\end{theorem}

\begin{proof}
The genus period ratio
\[
 B_{D,\varepsilon}=A_\varepsilon^6
              u_e^{-12h(e)\lambda_-}
\]
belongs to $K$. Indeed,
\[
 12h(e)\lambda_-=
 \frac{24h(e)h(D_-)}{\omega(D_-)}\in\mathbb Z,
\]
because $\omega(D_-)$ is $2$, $4$, or $6$.
Equations~\eqref{cmgenus:eq:four} and \eqref{cmgenus:eq:six} therefore
give $R_4(D)^3\in K$ and $R_6(D)^2\in K$.
The seed ratios are nonzero: $E_4$ vanishes only at the elliptic orbit
of discriminant $-3$, and $E_6$ only at that of discriminant $-4$,
whereas $D<-4$ is fundamental. Adjoining the two positive seed ratios
has degree at most $3\cdot2=6$ over $K$.

To prove that this one field suffices, use the reduced modular
polynomial of \eqref{newmodular:eq:reduced}:
\begin{equation}
 E_w=E_4^aE_6^b\Delta^m P_w(j),\qquad
 a\in\{0,1,2\},\quad b\in\{0,1\},\quad
 w=4a+6b+12m,
 \label{cmgenus:eq:smallpolynomial}
\end{equation}
where $m\geq0$ and $P_w\in\mathbb Q[X]$ is monic of degree $m$.
The algebraic reason is short: the weight condition forces every
monomial of $E_w\in\mathbb Q[E_4,E_6]$ to have exponents
$a+3r,b+2s$ with $r+s=m$. Substituting
$E_4^3=\Delta j$ and $E_6^2=\Delta(j-1728)$ gives the displayed
form; the constant term $E_w=1+O(q)$ makes $P_w$ monic.

The signed absolute product gives the formula without fractional
root extraction
\begin{equation}
 R_w(D)=R_4(D)^aR_6(D)^b B_{D,\varepsilon}^{m}
       \left|\frac{\operatorname{Res}(H_D^+,P_w)}
                    {\operatorname{Res}(H_D^-,P_w)}\right|.
 \label{cmgenus:eq:noroot}
\end{equation}
The resultant ratio is in $K$, because both genus factors have
coefficients there. Its absolute value is also in $K$, which is a
real field. Since the seed forms and $\Delta$ never vanish at these
CM points, the denominator condition on $R_w$ guarantees a nonzero
denominator resultant. This proves the field assertion.
\end{proof}

\begin{corollary}[Quadratic closure for the recorded discriminants]
\label{cmgenus:cor:closure}
For $D=-15,-20,-39$, every defined genus ratio $R_w(D)$ with even
$w\geq4$ lies in $\mathbb Q(\sqrt5)$, $\mathbb Q(\sqrt5)$, or
$\mathbb Q(\sqrt{13})$, respectively. Each is computable in that
quadratic field using the reduced polynomial $P_w$, whose degree
is at most $\lfloor w/12\rfloor$.
\end{corollary}
\begin{proof}
Theorems~\ref{cmgenus:thm:four} and \ref{cmgenus:thm:six} place both
seed ratios in the indicated quadratic field. Formula
\eqref{cmgenus:eq:noroot} completes the proof.
\end{proof}

For example, $E_{14}=E_4^2E_6$ gives
$R_{14}(D)=R_4(D)^2R_6(D)$.
At weight $16$, $a=m=1,b=0$ and
$P_{16}(X)=X-3456000/3617$, so
\[
 R_{16}(D)=R_4(D)B_{D,\varepsilon}
    \left|\frac{H_D^+(3456000/3617)}
                 {H_D^-(3456000/3617)}\right|.
\]
The equal degrees of the two genus factors cancel the signs in the
linear-polynomial resultants. This formula, together with the three
explicit pairs $H_D^\pm$, generates further exact quadratic identities
using rational arithmetic alone.
